% ROZDZIAŁ 27 -- ujęcie miarowe; definicje geometryczne w osobnym skrypcie.
\chapter[Miary na rozmaitościach]{Miary na rozmaitościach: ujęcie miarowo-teoretyczne}
\label{ch:miary-rozmaitosci-27}

\section{Zakres rozdziału i ustalone obiekty}
W opracowaniu \emph{Geometria i topologia różniczkowa}, w rozdziale
\emph{Tensor metryczny, długość i miara}, skonstruowano już
miarę riemannowską, wyprowadzono jej niezależność od mapy,
wprowadzono metrykę indukowaną i obliczono całki po gładkich
podrozmaitościach. W rozdziale o formach udowodniono twierdzenie
Stokesa. Nie powtarzamy tych dowodów ani przykładów
sfery i paraboloidy. Zamiast tego badamy miarę riemannowską
jako obiekt teorii miary: regularność, zbiory zerowe,
gęstości Radona--Nikodyma, transport miar i produkty.

Przez $(M^n,g)$ rozumiemy w tym rozdziale niepustą, gładką,
Hausdorffowską rozmaitość bez brzegu, spełniającą drugi aksjomat przeliczalności,
wymiaru $n\ge1$, wyposażoną w gładką metrykę riemannowską.
Nie wymagamy orientowalności. Symbol $\mathcal B(M)$ oznacza
$\sigma$-algebrę borelowską, a $\mu_g$ oznacza miarę
borelowską określoną lokalnie przez
\begin{equation}\label{eq:g-density-27}
  \mu_g(A\cap U)=\int_{x(A\cap U)}
  \sqrt{\det G_x(u)}\,\mathrm d\lambda_n(u),
\end{equation}
gdzie $(U,x)$ jest mapą, $A\in\mathcal B(M)$,
a $G_x(u)=(g_{ij}(x^{-1}(u)))$ jest macierzą metryki
w bazie współrzędnościowej w punkcie $x^{-1}(u)$.
Jest to przypomnienie definicji, nie nowa konstrukcja.
Jej uzupełnienie będziemy oznaczać przez
$(\overline{\mathcal B(M)}^{\,\mu_g},\overline\mu_g)$.

\section{Regularność i $\sigma$-skończoność}
\begin{twierdzenie}[Miara objętościowa jest miarą Radona]
\label{thm:riemann-radon-27}
Niech $(M^n,g)$ będzie gładką rozmaitością riemannowską
bez brzegu, Hausdorffa, spełniającą drugi aksjomat przeliczalności,
wymiaru $n\ge1$. Wówczas
$\mu_g$ jest lokalnie skończoną, $\sigma$-skończoną
miarą Radona na $\mathcal B(M)$. W szczególności
$\mu_g(K)<\infty$ dla każdego zwartego $K\subseteq M$.
\end{twierdzenie}
\begin{proof}
Weźmy punkt $p\in M$ i mapę $(U,x)$ wokół $p$.
Wybieramy kulę euklidesową $B$ tak, aby
$\overline B\subset x(U)$, a jej przeciwobraz $K=x^{-1}(\overline B)$
był zwartym otoczeniem punktu $p$. Gęstość
$\rho(u)=\sqrt{\det G_x(u)}$ jest dodatnia
i ciągła, zatem na $\overline B$ jest ograniczona
z góry i z dołu przez dodatnie stałe. Z
\eqref{eq:g-density-27} wynika $\mu_g(K)<\infty$.
Dzięki zwartości dowolny zwarty $L\subset M$ pokrywamy
skończenie wieloma takimi otoczeniami, skąd
$\mu_g(L)<\infty$.

Rozmaitość $M$ jest lokalnie zwarta i ma przeliczalną
bazę, więc jest $\sigma$-zwarta: wybieramy
przeliczalnie wiele otoczeń o zwartych domknięciach,
które ją pokrywają. W konsekwencji $\mu_g$ jest
$\sigma$-skończona. Ponadto $M$ jest przestrzenią
metryzowalną i lokalnie zwartą. Stosujemy
udowodnione wcześniej twierdzenie o regularności
lokalnie skończonych miar borelowskich na takich
przestrzeniach (Rozdział~\ref{ch:radon-measures-20}),
aby otrzymać regularność wewnętrzną przez zbiory
zwarte i zewnętrzną przez zbiory otwarte.
\end{proof}

\begin{uwaga}[Miara borelowska a jej uzupełnienie]
Nie należy utożsamiać $\mathcal B(M)$ z jej
uzupełnieniem. Podzbiory borelowskiego zbioru
miary zero mogą nie być borelowskie, ale należą do
$\overline{\mathcal B(M)}^{\,\mu_g}$.
Definicja miary Radona dotyczy w tym rozdziale
miary borelowskiej; do całkowania funkcji
mierzalnych względem jej uzupełnienia używamy
standardowego rozszerzenia całki.
\end{uwaga}

\section{Zbiory o zerowej objętości}
\begin{twierdzenie}[Lokalne kryterium zbiorów zerowych]
\label{thm:null-local-27}
Niech $(M^n,g)$ będzie gładką rozmaitością riemannowską
bez brzegu, Hausdorffa, spełniającą drugi aksjomat przeliczalności,
wymiaru $n\ge1$. Niech $A\in\mathcal B(M)$, a
$\{(U_j,x_j)\}_{j\ge1}$ będzie przeliczalnym atlasem
pokrywającym $M$. Wówczas
\[
 \mu_g(A)=0\quad\Longleftrightarrow\quad
 \lambda_n\bigl(x_j(A\cap U_j)\bigr)=0
 \quad\text{dla każdego }j\ge1.
\]
To samo kryterium działa dla zbiorów mierzalnych
w uzupełnionej $\sigma$-algebrze, gdy ich obrazy
rozumiemy jako zbiory mierzalne w sensie Lebesgue'a.
\end{twierdzenie}
\begin{proof}
Na każdej mapie gęstość $\rho_j$ jest dodatnia
i ciągła. Dla zbioru borelowskiego $E\subseteq x_j(U_j)$
mamy $\int_E\rho_j\,\mathrm d\lambda_n=0$
wtedy i tylko wtedy, gdy $\lambda_n(E)=0$:
wystarcza wyczerpać $x_j(U_j)$ zbiorami zwartymi,
na których $\rho_j$ ma dodatnie minimum.
Z \eqref{eq:g-density-27} otrzymujemy kryterium
lokalne, a przeliczalność atlasu pozwala
zsumować je na całym $M$. Jeżeli $A$ należy do uzupełnienia, istnieją zbiory borelowskie
$B,Z\subseteq M$ takie, że $A\triangle B\subseteq Z$ i $\mu_g(Z)=0$.
W każdej mapie obraz różnicy symetrycznej jest zawarty w
$x_j(Z\cap U_j)$, który na mocy części borelowskiej ma miarę
Lebesgue'a zero. Zupełność miary Lebesgue'a zapewnia zatem
mierzalność $x_j(A\cap U_j)$ oraz równość jego miary z miarą
$x_j(B\cap U_j)$. Ponadto $\overline\mu_g(A)=\mu_g(B)$.
Stosując część borelowską do $B$, otrzymujemy obie implikacje
dla uzupełnienia.
\end{proof}

\begin{wniosek}[Obrazy immersji niższego wymiaru]
\label{cor:submanifold-null-27}
Niech $(M^n,g)$ będzie gładką rozmaitością riemannowską
bez brzegu, Hausdorffa, spełniającą drugi aksjomat przeliczalności,
wymiaru $n\ge1$. Niech $N^k$ będzie gładką rozmaitością bez brzegu,
Hausdorffa, spełniającą drugi aksjomat przeliczalności, $0\le k<n$,
a $F:N\to M$ immersją klasy $C^1$. Wówczas $S=F(N)$
jest zbiorem borelowskim w $M$ i $\mu_g(S)=0$.
W szczególności teza zachodzi dla każdej osadzonej
podrozmaitości gładkiej wymiaru $k<n$.
\end{wniosek}
\begin{proof}
Dla $k=0$ przestrzeń $N$ jest dyskretna i przeliczalna,
więc jej obraz jest przeliczalną sumą punktów o zerowej
$\mu_g$-mierze. Niech zatem $k\ge1$.
Ustalmy $p\in N$ i gładkie mapy źródła i celu wokół $p$
oraz $F(p)$. Po ograniczeniu dziedziny mapy źródła tak, aby jej obraz
przez $F$ mieścił się w dziedzinie mapy celu, otrzymujemy
współrzędnościowe przedstawienie $F$, którego macierz
pochodnej w punkcie odpowiadającym $p$ ma rząd $k$. Po ewentualnej permutacji współrzędnych celu
pierwsze $k$ składowych tworzy odwzorowanie $P$ o odwracalnej
pochodnej w punkcie odpowiadającym $p$.
Z twierdzenia o odwzorowaniu odwrotnym, po ograniczeniu
dziedziny do otoczenia $V_p$, odwzorowanie $P$ jest
dyfeomorfizmem klasy $C^1$ na otwarty zbiór $W\subseteq\mathbb R^k$.
Pozostałe składowe, złożone z $P^{-1}$, tworzą funkcję
ciągłą $\psi:W\to\mathbb R^{n-k}$. Obraz $F(V_p)$
w wybranej gładkiej mapie celu jest więc wykresem $\psi$.
Wykres jest domknięty względem $W\times\mathbb R^{n-k}$,
zatem borelowski w $\mathbb R^n$. Każdy jego przekrój
pionowy nad punktem $W$ jest singletonem, a poza $W$ pusty.
Ponieważ $n-k\ge1$, twierdzenie Tonellego daje zerową
$n$-wymiarową miarę Lebesgue'a wykresu. Wzór
\eqref{eq:g-density-27} daje $\mu_g(F(V_p))=0$.
Takie uzasadnienie nie wymaga użycia map klasy $C^1$
jako map gładkiego atlasu celu.
Ponieważ $N$ ma przeliczalną bazę, z pokrycia
$(V_p)_{p\in N}$ wybieramy podpokrycie
przeliczalne $(V_j)$. Zatem
$S=\bigcup_jF(V_j)$. Każdy $F(V_j)$ jest borelowski
w dziedzinie gładkiej mapy celu, a więc także w $M$,
ponieważ dziedzina mapy jest otwarta. Stąd $S$ jest
borelowski i z przeliczalnej podaddytywności
$\mu_g(S)\le\sum_j\mu_g(F(V_j))=0$.
Argument nie zakłada, że \emph{cały} obraz immersji
jest lokalnie osadzoną podrozmaitością w otoczeniu
każdego punktu obrazu: przy samoprzecięciach lub
nagromadzeniach może to być nieprawdą.
\end{proof}

\begin{uwaga}[Dwie różne miary na tym samym zbiorze]
Wniosek nie oznacza, że długość krzywej lub pole
powierzchni znika. Dla osadzonej $k$-wymiarowej
podrozmaitości $S$, miara $\mu_{g|_S}$ może być
dodatnia, choć $\mu_g(S)=0$, gdy $k<n$.
To porównanie prowadzi do miar Hausdorffa;
identyfikację $\mu_{g|_S}$ z odpowiednio
unormowaną $\mathcal H^k$ uzasadnimy
w Rozdziale~\ref{ch:area-coarea-28}.
\end{uwaga}

\section{Miary z gęstością i rozkład Lebesgue'a}
\begin{twierdzenie}[Miary bezwzględnie ciągłe względem objętości]
\label{thm:rn-manifold-27}
Niech $(M^n,g)$ będzie gładką rozmaitością riemannowską
bez brzegu, Hausdorffa, spełniającą drugi aksjomat przeliczalności,
wymiaru $n\ge1$. Niech
$\nu$ będzie dodatnią, $\sigma$-skończoną
miarą na $\mathcal B(M)$, przy czym
$\nu\ll\mu_g$. Wtedy istnieje nieujemna funkcja
borelowska $h$, jednoznaczna $\mu_g$-prawie
wszędzie, taka że dla każdego $A\in\mathcal B(M)$
\[
 \nu(A)=\int_A h(p)\,\mathrm d\mu_g(p).
\]
Jeżeli dodatkowo $\nu$ jest lokalnie skończona,
to $h\in L^1_{\mathrm{loc}}(M,\mu_g)$,
tj. $\int_K h\,\mathrm d\mu_g<\infty$
dla każdego zwartego $K\subseteq M$.
\end{twierdzenie}
\begin{proof}
Miara $\mu_g$ jest $\sigma$-skończona na mocy
Twierdzenia~\ref{thm:riemann-radon-27}.
Stosujemy twierdzenie Radona--Nikodyma
(Rozdział~\ref{ch:radon-nikodym-18})
do pary $(\nu,\mu_g)$. Otrzymujemy nieujemną funkcję $h$ mierzalną względem $\mathcal B(M)$,
ponieważ twierdzenie stosujemy na tej właśnie $\sigma$-algebrze. Jeżeli $\nu$ jest lokalnie
skończona, każdy zwarty $K$ można przykryć
skończenie wieloma otoczeniami o skończonej
$\nu$-mierze, więc $\nu(K)<\infty$.
Wstawienie $A=K$ daje lokalną całkowalność.
\end{proof}

\begin{przyklad}[Gęstość w mapie]
Niech $(U,x)$ będzie mapą rozmaitości riemannowskiej
$(M^n,g)$, $V=x(U)\subseteq\mathbb R^n$, a $\nu$ dodatnią
$\sigma$-skończoną miarą borelowską na $M$.
Załóżmy, że miara obrazowa $x_*(\nu|_U)$ ma gęstość
$q:V\to[0,+\infty]$ względem $\lambda_n|_V$.
Wtedy gęstość $\nu|_U$ względem $\mu_g|_U$ wynosi
$h(p)=q(x(p))/\sqrt{\det G_x(x(p))}$ dla
$\mu_g$-prawie wszystkich $p\in U$. Istotnie, dla każdego
borelowskiego $A\subseteq U$ z lokalnego wzoru na miarę
riemannowską i definicji miary obrazowej otrzymujemy
\[
 \int_A h\,\mathrm d\mu_g
 =\int_{x(A)}q(a)\,\mathrm d\lambda_n(a)=\nu(A).
\]
Dodatniość gęstości riemannowskiej uzasadnia dzielenie.
Wyrażenie $q$ zmienia się po zmianie mapy;
$h$ jest funkcją skalarną na rozmaitości.
\end{przyklad}

\begin{twierdzenie}[Rozkład na część objętościową i osobliwą]
\label{thm:leb-manifold-27}
Niech $(M^n,g)$ będzie gładką rozmaitością riemannowską
bez brzegu, Hausdorffa, spełniającą drugi aksjomat przeliczalności,
wymiaru $n\ge1$. Jeżeli $\nu$ jest dodatnią, $\sigma$-skończoną
miarą na $\mathcal B(M)$, to ma jednoznaczny rozkład
$\nu=h\mu_g+\nu_{\mathrm s}$, gdzie
$h\ge0$ jest borelowska, a
$\nu_{\mathrm s}\perp\mu_g$. Jednoznaczność $h$ rozumiemy
z dokładnością do równości $\mu_g$-prawie wszędzie.
\end{twierdzenie}
\begin{proof}
Jest to zastosowanie twierdzenia o rozkładzie
Lebesgue'a (Rozdział~\ref{ch:ac-singular-17})
i twierdzenia Radona--Nikodyma do części
bezwzględnie ciągłej. Jednoznaczność obu składników
wynika z tych twierdzeń.
\end{proof}

\section{Przenoszenie objętości przez dyfeomorfizm}
\begin{definicja}[Jakobian objętościowy]
\label{def:volume-jacobian-27}
Niech $(M^n,g)$ i $(N^n,h)$ będą rozmaitościami
riemannowskimi wymiaru $n\ge1$, gładkimi, bez brzegu,
Hausdorffa i spełniającymi drugi aksjomat przeliczalności,
a $F:M\to N$ dyfeomorfizmem klasy $C^1$.
Dla $p\in M$ wybieramy ortonormalne bazy
w $T_pM$ i $T_{F(p)}N$, a przez $A_p$ oznaczamy
macierz $\mathrm dF_p$ w tych bazach. Definiujemy
$J_{g,h}F(p)=|\det A_p|$. Zmiana każdej z baz mnoży macierz z odpowiedniej strony
przez macierz ortogonalną, której wyznacznik ma moduł jeden.
Wartość nie zależy więc od wyboru baz.
Niech $u,v$ będą mapami wokół $p,F(p)$, a $G=G_u$ i $H=H_v$
macierzami metryk w tych współrzędnych. Wówczas
\[
 J_{g,h}F(p)=|\det D(v\circ F\circ u^{-1})(u(p))|
 \frac{\sqrt{\det H(v(F(p)))}}{\sqrt{\det G(u(p))}}.
\]
\end{definicja}

\begin{twierdzenie}[Zamiana zmiennych na rozmaitościach]
\label{thm:manifold-change-27}
Niech $(M^n,g)$ i $(N^n,h)$ będą gładkimi rozmaitościami
riemannowskimi bez brzegu, Hausdorffa, spełniającymi drugi
aksjomat przeliczalności, wymiaru $n\ge1$.
Niech $F:M\to N$ będzie dyfeomorfizmem klasy $C^1$,
a $J_{g,h}F$ jego jakobianem objętościowym.
Dla każdej funkcji borelowskiej $f:N\to[0,+\infty]$ zachodzi
\begin{equation}\label{eq:manifold-change-27}
 \int_N f(q)\,\mathrm d\mu_h(q)=
 \int_M f(F(p))J_{g,h}F(p)\,\mathrm d\mu_g(p).
\end{equation}
Ta sama równość zachodzi dla funkcji ze znakiem,
gdy $f\in L^1(N,\mu_h)$.
W szczególności $(F_*\mu_g)\ll\mu_h$ i
\[
 \frac{\mathrm d(F_*\mu_g)}{\mathrm d\mu_h}(q)
 =\frac{1}{J_{g,h}F(F^{-1}(q))}
 \quad\text{dla }\mu_h\text{-prawie wszystkich }q.
\]
\end{twierdzenie}
\begin{proof}
Uzasadnijmy najpierw wzór współrzędnościowy z definicji
jakobianu. Ustalmy $p\in M$ i mapy $u,v$ wokół $p,F(p)$,
ograniczając w razie potrzeby dziedzinę $u$ tak, aby jej obraz
przez $F$ mieścił się w dziedzinie $v$.
Niech $T=v\circ F\circ u^{-1}$, $D=DT(u(p))$,
$G=G_u(u(p))$ oraz $H=H_v(v(F(p)))$.
Niech kolumny macierzy $P$ i $Q$ będą współrzędnymi
wektorów wybranych baz ortonormalnych względem baz
współrzędnościowych źródła i celu. Wtedy
$P^{\mathsf T}GP=I$ oraz $Q^{\mathsf T}HQ=I$, a macierz
różniczki w bazach ortonormalnych to $A_p=Q^{-1}DP$.
Z równości wyznaczników dostajemy
$|\det P|=(\det G)^{-1/2}$ i
$|\det Q|=(\det H)^{-1/2}$, skąd
\[
 |\det A_p|=|\det D|\frac{\sqrt{\det H}}{\sqrt{\det G}}.
\]
W szczególności $J_{g,h}F$ jest ciągły i ściśle dodatni.
Dla różniczki odwrotnej macierz w tych samych bazach,
zamienionych rolami, to $A_p^{-1}$, zatem
$J_{h,g}(F^{-1})(F(p))=1/J_{g,h}F(p)$.

Najpierw załóżmy, że $f$ znika poza otwartym zbiorem $V$
zawartym w dziedzinie mapy celu, a $F^{-1}(V)$
mieści się w dziedzinie mapy źródła. Z definicji
$\mu_g,\mu_h$ oraz twierdzenia o zamianie
zmiennych w $\mathbb R^n$ otrzymujemy
\eqref{eq:manifold-change-27}, po skróceniu
wyznaczników macierzy metryk: jeśli $T=v\circ F\circ u^{-1}$,
to iloczyn $J_{g,h}F(u^{-1}(a))\sqrt{\det G(a)}$
jest równy $|\det DT(a)|\sqrt{\det H(T(a))}$.
Jest to dokładnie gęstość występująca po euklidesowej zamianie zmiennych.

W ogólności dla każdego punktu $q\in N$ wybieramy otwarte
$V$ zawarte w dziedzinie mapy celu tak, aby $F^{-1}(V)$
mieścił się w dziedzinie mapy źródła. Istnienie takiego $V$
wynika z ciągłości $F^{-1}$. Z drugiego aksjomatu przeliczalności
wybieramy przeliczalne podpokrycie $(V_j)$. Tworzymy rozłączny
podział borelowski $B_1=V_1$, $B_j=V_j\setminus\bigcup_{i<j}V_i$.
Dla $f\chi_{B_j}$ stosujemy argument lokalny na $V_j$,
a następnie sumujemy równości, korzystając z twierdzenia
o zbieżności monotonicznej.

Aby obliczyć miarę obrazową, dla $A\in\mathcal B(N)$
stosujemy równanie do $F^{-1}$ i funkcji $\chi_{F^{-1}(A)}$ na $M$.
Ponieważ
$J_{h,g}(F^{-1})(q)=[J_{g,h}F(F^{-1}q)]^{-1}$,
dostajemy podaną pochodną Radona--Nikodyma.
Dla funkcji całkowalnych stosujemy równość
oddzielnie do części dodatniej i ujemnej.
\end{proof}

\begin{przyklad}[Konforemna zmiana metryki]
Niech $\varphi\in C^\infty(M;\mathbb R)$ i
$g'=e^{2\varphi}g$. Macierze spełniają
$G'=e^{2\varphi}G$, więc
$\det G'=e^{2n\varphi}\det G$. W konsekwencji
\[
 \mathrm d\mu_{g'}(p)=e^{n\varphi(p)}\,\mathrm d\mu_g(p).
\]
Nie oznacza to, że krzywizna obu metryk jest
taka sama: wzór dotyczy tylko miary objętościowej.
\end{przyklad}

\section{Miara na produkcie rozmaitości}
\begin{twierdzenie}[Miara produktu riemannowskiego]
\label{thm:product-riemann-27}
Niech $(M^n,g)$ i $(N^m,h)$ będą gładkimi rozmaitościami
riemannowskimi bez brzegu, Hausdorffa, spełniającymi drugi
aksjomat przeliczalności, wymiarów odpowiednio $n,m\ge1$. Na $M\times N$ rozważmy metrykę
$g\oplus h$. Wtedy na $\mathcal B(M\times N)$
zachodzi $\mu_{g\oplus h}=\mu_g\otimes\mu_h$.
Dla każdej funkcji borelowskiej $f:M\times N\to[0,+\infty]$ mamy
\[
 \int_{M\times N} f(p,q)\,\mathrm d\mu_{g\oplus h}(p,q)
 =\int_M\!\left(\int_N f(p,q)\,\mathrm d\mu_h(q)\right)
   \mathrm d\mu_g(p).
\]
\end{twierdzenie}
\begin{proof}
Ponieważ obie rozmaitości spełniają drugi aksjomat przeliczalności,
$\mathcal B(M\times N)=\mathcal B(M)\otimes\mathcal B(N)$:
każdy zbiór otwarty produktu jest przeliczalną sumą otwartych
prostokątów z baz obu czynników, a prostokąty borelowskie są
borelowskie dzięki ciągłości rzutów.
W mapach produktowych macierz metryki jest
blokowo diagonalna $\operatorname{diag}(G,H)$.
Stąd $\sqrt{\det(G\oplus H)}=
\sqrt{\det G}\sqrt{\det H}$. Dla zbiorów borelowskich $A\subseteq U$ i $B\subseteq V$,
gdzie $(U,u)$ i $(V,v)$ są mapami czynników, wzór na
gęstość i twierdzenie Tonellego dają
\[
 \mu_{g\oplus h}(A\times B)=
 \int_{u(A)}\int_{v(B)}
 \sqrt{\det G_u(a)}\sqrt{\det H_v(b)}
 \,\mathrm d\lambda_m(b)\,\mathrm d\lambda_n(a)
 =\mu_g(A)\mu_h(B).
\]
W iloczynie wartości miar stosujemy konwencję $0\cdot\infty=0$.
Obie miary ograniczone do $U\times V$ są więc równe
na prostokątach, a przez $\sigma$-skończoność czynników
i jednoznaczność miary produktowej są równe na całej
$\mathcal B(U\times V)$. Wybieramy przeliczalne atlasy
obu czynników i numerujemy wszystkie ich prostokąty
jako $(R_j)_{j\ge1}$. Dla dowolnego zbioru borelowskiego
$E\subseteq M\times N$ zbiory
$E_j=E\cap(R_j\setminus\bigcup_{i<j}R_i)$ są rozłączne,
borelowskie i pokrywają $E$. Równość miar na każdym
$E_j\subseteq R_j$ oraz przeliczalna addytywność
dają równość miar na $E$.
Ostatnia równość jest twierdzeniem Tonellego.
\end{proof}

\section{Od gładkich podrozmaitości do zbiorów prostowalnych}
W geometrii rozważaliśmy miarę indukowaną
na \emph{gładkich} podrozmaitościach. Obraz
odwzorowania Lipschitzowskiego nie musi być
gładką podrozmaitością, lecz nadal może
posiadać sensowną miarę $k$-wymiarową.
W następnym rozdziale pokażemy, jak
$\mathcal H^k$ mierzy takie obrazy i dlaczego
przy odwzorowaniach nieiniektywnych konieczny
jest czynnik krotności. Wprowadzone tam wzory
są ściśle miarowe: nie potrzebują globalnej
orientacji ani globalnych parametryzacji.
